% =============================================================
% Chapter 01 — Design Principles Before Writing Code
% Rewritten per Part 1 Deep Edit Report
% =============================================================

% ------------------------------------------------------------------
% Chapter Opening — Before you build: the failure modes
% ------------------------------------------------------------------

Before writing a single line of framework code, three failure modes deserve names.
The first is \textbf{API-centred architecture}: the codebase organises itself around
LLM request payloads rather than runtime responsibilities, so every change to a model
provider ripples through unrelated modules.
The second is \textbf{hidden state}: the agent's position in its own execution lives
in local variables or prompt text, making it invisible to tests and impossible to replay
deterministically.
The third is \textbf{undifferentiated failure}: a single \texttt{catch} block that
cannot tell a transient model outage from a malformed tool argument — both trigger the
same retry and produce the same log line.
These three failure modes are not hypothetical.
They appear in every production agent system that skipped the design step.
Part~1 is the antidote.

% ------------------------------------------------------------------
\section{Runtime Requirements}
\label{sec:runtime-requirements-recap}
% ------------------------------------------------------------------

Volume~1 answered the question \emph{what should an agent do}.
Volume~2 answers the question \emph{who owns what}.
The distinction matters because ownership is what allows each component to be
tested, replaced, and reasoned about in isolation.

Each requirement below is the direct inverse of a failure mode observed in
production agent systems.

\begin{enumerate}
  \item \textbf{Explicit state} — the run's position in the execution graph is
        always a named, persisted value.
  \item \textbf{Narrow contracts} — each component declares a single
        responsibility and depends only on the abstractions it needs.
  \item \textbf{Inward dependency} — infrastructure types (Redis keys, HTTP
        status codes, vendor SDK classes) never appear inside the runtime core.
  \item \textbf{Structured events} — every observable state transition emits a
        typed event that downstream subscribers can consume.
  \item \textbf{Typed failures} — \texttt{ModelFailure}, \texttt{ToolFailure},
        and \texttt{ValidationFailure} are distinct types that drive distinct
        recovery decisions.
  \item \textbf{Testability-first} — the full run loop can be exercised without
        a live model endpoint or a running persistence layer.
\end{enumerate}

\begin{table}[htbp]
\centering
\caption{Requirements mapped to the failure modes they prevent}
\label{tab:requirements-failure-modes}
\begin{tabular}{ll}
\toprule
\textbf{Requirement} & \textbf{Failure it prevents} \\
\midrule
Explicit state     & Hidden mutation — run position inferred from local variables or prompt text \\
Narrow contracts   & God-object runtime — planning, calling, executing, and emitting couple into one class \\
Inward dependency  & Infrastructure leakage — Redis types or vendor SDK classes appear in core logic \\
Structured events  & Silent black box — the only observable output is a final answer or a stack trace \\
Typed failures     & Undifferentiated catch — \texttt{ModelFailure} and \texttt{ToolFailure} trigger identical retry logic \\
Testability-first  & Slow test suite — every test requires a live model endpoint and a running Redis instance \\
\bottomrule
\end{tabular}
\end{table}

\begin{notebox}
  The six requirements are not stylistic preferences.
  They are constraints derived from failure modes in shipping systems.
  Treat any design decision that violates one of them as a debt item,
  not an acceptable trade-off.
\end{notebox}

\begin{takeawaybox}
  Volume~2 is a book about ownership.
  Before writing a class, ask: which requirement does this class serve,
  and what failure mode does that requirement prevent?
\end{takeawaybox}

% ------------------------------------------------------------------
% Figure 1.1 — Framework layer architecture and dependency direction
% ------------------------------------------------------------------
\begin{figure}[htbp]
\centering
\begin{tikzpicture}[
  layer/.style={draw, rounded corners=4pt, minimum width=9cm,
                minimum height=0.72cm, text centered, font=\small\ttfamily},
  proto/.style={layer,   fill=gray!15},
  sec/.style={layer,     fill=orange!18},
  obs/.style={layer,     fill=yellow!22},
  domlayer/.style={layer, fill=teal!18},
  orch/.style={layer,    fill=blue!15},
  core/.style={layer,    fill=blue!30},
  arr/.style={->, thick, >=stealth}
]
  \node[proto]    (proto)  at (0, 5.4) {protocols — HTTP / MCP / A2A adapters};
  \node[sec]      (sec)    at (0, 4.4) {security — PolicyEngine, SecretResolver};
  \node[obs]      (obs)    at (0, 3.4) {observability — EventPublisher, exporters};
  \node[domlayer] (domain) at (0, 2.4) {tooling + memory — ToolRegistry, MemoryStore};
  \node[orch]     (orchn)  at (0, 1.4) {orchestrator — AgentOrchestrator, AgentRunner};
  \node[core]     (coren)  at (0, 0.4) {runtime/run — Run, RunState, RunRepository};

  \draw[arr] (proto.south)  -- (sec.north);
  \draw[arr] (sec.south)    -- (obs.north);
  \draw[arr] (obs.south)    -- (domain.north);
  \draw[arr] (domain.south) -- (orchn.north);
  \draw[arr] (orchn.south)  -- (coren.north);

  \node[font=\footnotesize, right=0.4cm of coren]  {Ch\,1--2};
  \node[font=\footnotesize, right=0.4cm of orchn]  {Ch\,2--3};
  \node[font=\footnotesize, right=0.4cm of domain] {Ch\,3};
  \node[font=\footnotesize, right=0.4cm of obs]    {Ch\,2,\,4};
  \node[font=\footnotesize, right=0.4cm of sec]    {Ch\,4};
  \node[font=\footnotesize, right=0.4cm of proto]  {Ch\,4};

  \node[font=\small\itshape, below=0.3cm of coren]
    {Dependencies point inward — upper layers depend on lower; never the reverse.};
\end{tikzpicture}
\caption{Framework layer architecture. Arrows show the allowed dependency
  direction. Each layer may only import types from layers below it.
  Chapter references indicate where each layer is implemented.}
\label{fig:layer-architecture}
\end{figure}

% ------------------------------------------------------------------
\section{Testability and Observability}
\label{sec:testability-and-observability}
% ------------------------------------------------------------------

Testability and observability are usually treated as engineering virtues added after
a system is working.
In agent frameworks they must be structural from day one, because the cost of retrofit
is disproportionate: adding structured events after the fact requires touching every
state transition, and adding dependency seams after the fact requires rewriting every
constructor.

\subsection{Architectural Principles}
\label{subsec:architectural-principles}

\begin{table}[htbp]
\centering
\caption{Six architectural principles that make a runtime testable and observable}
\label{tab:arch-principles}
\begin{tabular}{lll}
\toprule
\textbf{Principle} & \textbf{Mechanism} & \textbf{What it enables} \\
\midrule
Seam at every boundary    & Interfaces, not classes, in constructor parameters & Fake any dependency in tests \\
Events, not logs          & Typed event objects on a publish–subscribe bus     & Machine-readable traces \\
State is data             & Serialisable run state, not thread state           & Deterministic replay \\
Failure is typed          & Sealed failure hierarchy                           & Per-failure-type retry logic \\
Side effects at the edge  & I/O only in adapters and repositories              & Pure-function core loop \\
Builder validates         & Required deps checked at construction, not first use & Fast, clear test failures \\
\bottomrule
\end{tabular}
\end{table}

\begin{warningbox}
  Retrofit observability is expensive.
  If you add structured events after the state machine is written, you will touch
  every \texttt{transition()} call site.
  Wire the \texttt{EventPublisher} interface into the runtime constructor on day one,
  even if the first implementation does nothing but discard events.
\end{warningbox}

\begin{tipbox}
  The cost of a fake dependency is close to zero.
  A \texttt{FakeModelClient} that returns a fixed response is twenty lines.
  The benefit is a test suite that runs in milliseconds with no network access.
  Write the fake before you write the real implementation.
\end{tipbox}

\subsection{Three-Level Observability}
\label{subsec:observability-hierarchy}

The runtime produces observable evidence at three levels.
At the \emph{run} level, a \texttt{RunSnapshot} records the run's identity, current
state, and timestamps.
At the \emph{step} level, a \texttt{RunStep} records which state the agent was in
when each unit of work began and ended.
At the \emph{dependency} level, individual collaborators (model client, tool registry,
memory store) record their inputs and outputs independently of the run loop.
The three levels compose: a single externally visible failure can be traced from the
run's terminal state back through the failing step to the dependency call that produced
the typed failure.

\subsection{Where the Evidence Lives}
\label{subsec:what-to-make-observable}

The step record lives in \texttt{RunStep} (\path{runtime/run/RunStep.java})\footnote{\texttt{RunStep.java}:
\url{https://github.com/tnhmar/agent-framework/blob/main/src/main/java/com/agentruntime/runtime/run/RunStep.java}}.
The durable run state lives in \texttt{RunSnapshot} (\path{runtime/run/RunSnapshot.java})\footnote{\texttt{RunSnapshot.java}:
\url{https://github.com/tnhmar/agent-framework/blob/main/src/main/java/com/agentruntime/runtime/run/RunSnapshot.java}}.
The repository that writes both is \texttt{RunRepository} (\path{runtime/run/RunRepository.java})\footnote{\texttt{RunRepository.java}:
\url{https://github.com/tnhmar/agent-framework/blob/main/src/main/java/com/agentruntime/runtime/run/RunRepository.java}}.
The three-part evidence trail described above is not a design aspiration.
It is the behaviour of the framework as it exists in the repository.

\begin{takeawaybox}
  Observability is not instrumentation added to a finished system.
  It is the shape of the system's data model.
  A run that records \texttt{RunStep} entries as it advances is observable by
  construction — no logging framework required.
\end{takeawaybox}

% ------------------------------------------------------------------
\section{Java Concurrency Models}
\label{sec:java-concurrency-models}
% ------------------------------------------------------------------

An agent runtime faces a concurrency choice early: should individual runs advance
concurrently, and if so, by what mechanism?
The options — platform threads, virtual threads, reactive streams, and structured
concurrency — differ in how they handle blocking calls, how they expose cancellation,
and how they interact with persistence.

\begin{table}[htbp]
\centering
\caption{Concurrency models compared for agent run execution}
\label{tab:concurrency-models}
\begin{tabular}{llll}
\toprule
\textbf{Model} & \textbf{Blocking cost} & \textbf{Cancellation} & \textbf{State replay} \\
\midrule
Platform threads        & High (1 MB stack each)  & Thread.interrupt()    & Possible \\
Virtual threads (JDK 21)& Low (continuation-based)& Thread.interrupt()    & Possible \\
Reactive streams        & None                    & Subscription.cancel() & Complex  \\
Structured concurrency  & Low                     & ShutdownOnFailure     & Possible \\
\bottomrule
\end{tabular}
\end{table}

\subsubsection{Why Sequential Within a Run}
\label{subsec:why-sequential}

The framework advances each run \emph{sequentially}: one step at a time, one state
transition per loop iteration.
This is a deliberate choice.
A run that advances two steps concurrently cannot guarantee that the step records
are written in causal order, which breaks deterministic replay.
Concurrency happens \emph{across} runs (many runs advancing simultaneously), not
\emph{within} a run.
This distinction is what allows the framework to offer cooperative pause and cancel
semantics: the pause signal is checked at state transition boundaries, where the
system is always in a consistent, persisted state.

\subsubsection{Framework Orchestration Types}
\label{subsec:concurrency-boundary-types}

The framework separates scheduling intent from execution mechanics.
\texttt{AgentOrchestrator} (\path{orchestrator/AgentOrchestrator.java})\footnote{\texttt{AgentOrchestrator.java}:
\url{https://github.com/tnhmar/agent-framework/blob/main/src/main/java/com/agentruntime/orchestrator/AgentOrchestrator.java}}
is the submission surface — it accepts work and returns a handle to the running agent.
\texttt{AgentRunner} (\path{orchestrator/AgentRunner.java})\footnote{\texttt{AgentRunner.java}:
\url{https://github.com/tnhmar/agent-framework/blob/main/src/main/java/com/agentruntime/orchestrator/AgentRunner.java}}
advances individual runs by submitting \texttt{AgentTask} units to a virtual-thread
executor.
\texttt{AgentLifecycle} (\path{orchestrator/AgentLifecycle.java})\footnote{\texttt{AgentLifecycle.java}:
\url{https://github.com/tnhmar/agent-framework/blob/main/src/main/java/com/agentruntime/orchestrator/AgentLifecycle.java}}
exposes the pause, resume, and cancel controls that Chapter~2 implements.

\begin{takeawaybox}
  Do not make the concurrency model visible inside the run loop.
  The loop transitions states; the scheduler decides when to advance.
  Keeping those concerns separate is what allows the pause-and-resume contract
  to be honoured without cooperative thread interruption.
\end{takeawaybox}

% ------------------------------------------------------------------
\section{Dependency Injection and Modularity}
\label{sec:dependency-injection-and-modularity}
% ------------------------------------------------------------------

\subsection{Interface Catalogue}
\label{subsec:interface-catalogue}

The table below is the canonical contract map for Part~1.
Every name listed here exists in the framework repository.
No name not in this table may appear in a code listing or prose description in
Chapters~1 through~4.

\begin{table}[htbp]
\centering
\caption{Framework interface catalogue — Part 1 reference}
\label{tab:interface-catalogue}
\scriptsize
\begin{tabular}{lllll}
\toprule
\textbf{Interface / Type} & \textbf{Responsibility} & \textbf{Layer} & \textbf{Package} & \textbf{Chapter} \\
\midrule
\texttt{AgentOrchestrator} & Own a run from submission to termination     & Runtime core  & \texttt{orchestrator}    & Ch~2 \\
\texttt{AgentRunner}       & Schedule and advance individual runs         & Runtime core  & \texttt{orchestrator}    & Ch~2 \\
\texttt{AgentLifecycle}    & Expose pause, resume, cancel controls        & Runtime core  & \texttt{orchestrator}    & Ch~2 \\
\texttt{RunRepository}     & Persist \texttt{Run} aggregate and step history & Runtime core & \texttt{runtime/run}   & Ch~2 \\
\texttt{RunState}          & Enumerate legal execution states             & Runtime core  & \texttt{runtime/run}     & Ch~1 \\
\texttt{RunSnapshot}       & Immutable durable projection of a run        & Runtime core  & \texttt{runtime/run}     & Ch~2 \\
\texttt{ModelClient}       & Provider-neutral language model invocation   & Model layer   & \texttt{modelclient}     & Ch~2 \\
\texttt{AgentDelegate}     & Delegate work to a sub-agent                 & Orchestration & \texttt{orchestrator}    & Ch~3 \\
\texttt{AgentTask}         & Represent a unit of scheduled work           & Orchestration & \texttt{orchestrator}    & Ch~2 \\
\texttt{ToolRegistry}      & Resolve tools by name, expose descriptors    & Tool layer    & \texttt{tooling}         & Ch~3 \\
\texttt{MemoryStore}       & Load and append memory records               & Memory layer  & \texttt{memory}          & Ch~3 \\
\texttt{MemoryPolicy}      & Window selection and retention decisions     & Memory layer  & \texttt{memory}          & Ch~3 \\
\texttt{EventPublisher}    & Structured runtime event emission            & Observability & \texttt{observability}   & Ch~2 \\
\texttt{PolicyEngine}      & Evaluate security policy before operations   & Security      & \texttt{security}        & Ch~4 \\
\texttt{SecretResolver}    & Resolve credentials without raw exposure     & Security      & \texttt{security}        & Ch~4 \\
\texttt{ProtocolAdapter}   & Translate protocol messages to commands      & Protocol edge & \texttt{protocols}       & Ch~4 \\
\bottomrule
\end{tabular}
\end{table}

\subsection{Core State Model}
\label{subsec:core-state-model}

The \texttt{RunState} enum is the vocabulary of the entire framework.
Every component that cares about what an agent is doing — the loop, the repository,
the event publisher, the interrupt handler — speaks in \texttt{RunState} values.
The \texttt{PAUSED} state is not a terminal state: it allows the run to be
resumed from exactly the position at which it was paused, with no work repeated.
The terminal states (\texttt{COMPLETED}, \texttt{FAILED}, \texttt{CANCELLED},
\texttt{TIMED\_OUT}) are sink states — once entered, no transition out is legal.

The listing below is the authoritative definition, taken ver